\section{Introduction}
\label{sec:intro}

Large language-model agents have moved from demonstration to deployment in scientific research itself: since the first demonstrations of autonomous experimentation and tool-augmented chemistry \cite{boiko2023emergent,bran2023chemcrow}, systems that plan and execute multi-step analyses now contribute to genome-scale model construction, autonomous chemistry, and end-to-end research pipelines \cite{qian2026towards,xie2026autonomous,kong2026autoresearch}, and surveys of scientific large language models document the breadth of this movement \cite{ai4science2023impact,zhang2024comprehensive,hogg2024machine}. This deployment rests on an architectural habit --- a single frontier policy model drives most or all reasoning steps --- even though the capability distribution across steps is anything but uniform. Steps that explore a dataset, check an intermediate result, or format an artifact rarely demand the same reasoning depth as steps that must synthesize a flawed plan into a working deliverable, and the efficiency literature already shows that smaller models outperform larger ones on many narrow, well-specified subtasks \cite{cao2026taskspecific,garg2025emissions}. Meanwhile the practical machinery of agent systems --- memory, skills, harnesses, orchestration layers --- has grown enormously more sophisticated \cite{zhou2026externalization,wei2026survey}. What has not grown is any principled answer to the allocation question this progress begs: \emph{which steps of a scientific agent actually need the frontier model, and how would one know?}

Existing evaluation practice cannot answer it. The dominant benchmarks score agents at the task level --- repository issue resolution, scientific workflows, equation discovery --- and report a single aggregate success rate \cite{yang2024sweagent,deng2025swebench,zhang2025swebench,sun2025scienceboard,shojaee2025llmsrbench,fang2025mlzero}. Task-level scores are silent about step-level structure: an agent at 40\% may be uniformly mediocre or excellent-except-one-step, and these two profiles demand opposite interventions. The tension is sharpened by the agentic self-improvement loop that dominates current practice --- reflection and self-refinement attach step-level critiques to a single policy \cite{yao2022react,shinn2023reflexion,madaan2023selfrefine}, and verifier-augmented training pipelines consume step-level signals wholesale \cite{pan2024training} --- so that a wrong picture of step structure propagates directly into training data. It is further sharpened by evidence that test-based benchmarks themselves can be gamed, which makes objective, hard-to-game measurement designs a first-order concern \cite{yu2026sweabs}. Where judges are needed to score open-ended steps, LLM-as-a-judge evaluation introduces its own reliability-validity problems \cite{gu2024survey,wei2024systematic,norman2026reliability}. A parallel literature on process supervision scores individual reasoning steps --- process reward models, step-level verification, and critiques of reward hacking show the appetite for step-level signals \cite{ma2024what,zhang2025groundedprm,pronesti2026beyond,zhang2026when,nie2026stop} --- but these scores are \emph{learned} judgments, typically distilled from models or humans, and inherit the circularity they are meant to resolve: the verifier is another model whose own competence boundary is unknown. Failure-attribution tools locate \emph{which} step failed post hoc \cite{deng2026defa}, yet attribution is not intervention: nothing in it establishes what would have happened had a different policy executed that step. To our knowledge, no existing protocol measures, by direct experiment, the causal consequence of substituting the model at one step while holding everything else fixed.

We close this gap with \textbf{step-level counterfactual replacement}. The protocol freezes a complete trajectory produced by a frontier policy model on a scientific-computing task, selects a step, and re-executes the trajectory from that step with one substitution --- the step's action is produced by a different model, from an identical prompt, in an identical sandbox, with identical remaining budget --- while all other steps retain the frontier policy. Success is judged only by objective executors (unit tests, file assertions, programmatic checks), never by model opinion, and every replacement arm is paired with a frontier-resampling arm that controls for stochasticity of re-execution itself. The design inherits the intervention logic of counterfactual evaluation while targeting the step, the natural unit at which agent frameworks are actually composed and at which model routing decisions must ultimately be made \cite{li2026progrouter,sun2026cborchestrator}. Because outcomes are objective, the protocol yields ground truth that process-reward approaches approximate and judge-based evaluation cannot certify --- a distinction of practical weight given evidence that LLM judges track human evaluation only imperfectly even in software engineering \cite{wang2025llms}, and given the provenance-auditing needs of agentic workflows \cite{souza2025agents}.

Applying the protocol to scientific data-analysis agents --- 767 controlled runs over 38 step instances across four replacement policies --- yields three findings. First, replacement damage is concentrated and heterogeneous: difference-in-differences against matched control steps attributes step-specific losses of 0.21 (mini), 0.25 (DeepSeek) and 0.71 (Qwen-14B) to six signal steps, and error-family coding shows the losses arise by different mechanisms --- protocol lapses for API models, a model-idiosyncratic error family for the open-weights model --- rather than by uniform capability shortfall. Second, dense scans of complete trajectories separate steps into healthy, capability-dependent and cliff classes, and interventions validate the taxonomy: structured feedback repairs a substantial fraction of failures (+0.21), suffix takeover --- splicing a cheaper edit and letting the frontier model continue --- amplifies damage, and doubling open-weights model size does not dissolve step-level failures. Third, and most constructively, seven verifiability features of a step --- coded blind to all replacement outcomes from frozen reading materials, with pre-registered adjudication and a frozen timestamp --- predict step-level replaceability under task-level cross-validation, meeting both pre-registered criteria ($\Delta$RMSE $\leq$ 0, $\Delta$$\rho$ $\geq$ 0.1). Steps lacking an objective correctness standard are reliably the safest to delegate to cheaper models; deliverable-producing steps are not.

Our contributions are:

\begin{itemize}
\item a \textbf{measurement protocol} --- step-level counterfactual replacement with objective adjudication and resampling controls --- that turns step-level model dependence into an experimentally measurable quantity rather than a judged one;
\item an \textbf{anatomy of frontier dependence} in scientific agents: concentration at identifiable steps, a three-class step taxonomy, and mechanism-level evidence that failures are discrete, model-characteristic events (protocol lapses, idiosyncratic error families) rather than graded attenuation;
\item \textbf{intervention results} that bound repair strategies: structured feedback helps, naive splicing hurts, and scale alone does not fix step failures;
\item a \textbf{pre-registered, blind-coded prediction test} showing that verifiability features carry independent signal about where cheaper models suffice, establishing verifiability as a measurable design principle for allocating frontier capability in agent pipelines.
\end{itemize}

Section~\ref{sec:results} reports the four result blocks; Methods specifies the protocol, the annotation and prediction designs, and all pre-registrations; the Discussion locates the findings against process supervision and model routing, and states the limits of a single-skeleton, single-domain, 38-point study.
